\section{Discussion}
\label{sec:discussion}

\para{What do the results mean?} Across both tasks, prompt-only frontier LLMs outperform the implemented classical baselines by large margins. The result is strongest in the \minimalprompt condition, which suggests that the models can infer short-horizon continuation cues from recent listening history or playlist context without additional user-summary text. In practical terms, a user who exports listening data and asks a frontier model to rank a fixed candidate set can obtain recommendations that are much stronger than lightweight popularity or cooccurrence heuristics.

\para{Why does the simple prompt work better?} The most plausible explanation is that once the candidate set is fixed, the task is mainly about local coherence. Recent tracks provide strong evidence about artist clusters, genre, mood, and sequence flow. Compact profile summaries may dilute that signal by introducing broader but less relevant preferences. The \mpd ablation is especially consistent with this interpretation because the profile prompt is reliably worse for both models.

\para{Why are the gains plausible?} The candidate sets are small, and the task rewards semantic matching between the recent context and one correct continuation. Frontier LLMs are likely strong at recognizing style, artist, and playlist-theme compatibility in text form. That capability can dominate lightweight history-based heuristics, especially when the context is coherent. The results therefore support a narrow but credible claim: semantic reasoning over recent history is enough to beat simple public baselines in a reranking setting.

\para{Popularity and diversity.} The LLM conditions are not merely replaying the popularity head. They consistently achieve higher mean top-10 popularity percentiles and slightly higher coverage than the classical baselines. This does not prove better novelty in a user-centric sense, but it does show that the gains are compatible with less head-heavy recommendation lists.

\para{Limitations.} The paper does not compare against Spotify's production recommender. We evaluate public datasets, a closed 20-item candidate pool, and two lightweight baselines. The \lastfm experiment required switching from a sparse local subset to a larger remote dataset slice, and the \mpd evaluation uses only one downloaded slice rather than the full benchmark. Harder negative sampling could reduce all scores. Finally, we perform no human evaluation, so the study measures offline ranking accuracy rather than subjective playlist quality or satisfaction.

\para{Broader implications.} The main broader implication is about portability rather than raw leaderboard performance. Prompt-only recommendation offers a user-controlled interface to personal history that does not require access to a platform's internal feature store or retraining pipeline. At the same time, our findings should not be overstated. A production recommender typically uses richer feedback signals, larger candidate generation systems, and long-term optimization targets that are absent here. The most realistic path forward may therefore be hybrid: use classical retrieval or sequential models for candidate generation, then use an LLM as a semantic reranker or explanation layer.
